\section{Introduction}
\label{sec:intro}

Bài toán N-Queens yêu cầu đặt $N$ quân hậu trên bàn cờ kích thước $N\times N$ sao cho không có hai quân hậu nằm trên cùng một hàng, cột hoặc đường chéo. Cùng một tập điều kiện này có thể được biểu diễn bằng các biến Boolean ứng với từng ô cờ, các biến nguyên xác định vị trí hậu trên từng hàng, hoặc một hệ ràng buộc tuyến tính. Sự đa dạng đó làm cho N-Queens trở thành một bài toán phù hợp để khảo sát ảnh hưởng của cách mô hình hóa đến hiệu năng của các bộ giải.

Trong cách tiếp cận SAT, mỗi ô cờ được biểu diễn bằng một biến Boolean và các điều kiện hợp lệ được chuyển thành công thức ở dạng chuẩn hội (CNF). Hai thành phần cơ bản là ràng buộc \emph{ít nhất một} (At Least One -- ALO) và \emph{nhiều nhất một} (At Most One -- AMO). Kết hợp ALO với AMO bảo đảm có đúng một quân hậu trên mỗi hàng và mỗi cột; AMO trên các đường chéo bảo đảm các quân hậu không tấn công nhau theo đường chéo. Tuy nhiên, một ràng buộc AMO có nhiều cách mã hóa tương đương về tập nghiệm trên các biến ban đầu, nhưng khác nhau về số biến phụ, số mệnh đề và cấu trúc suy diễn.

Theo yêu cầu Bài tập lớn số 1 của học phần Các vấn đề hiện đại KHMT \citep{assignment}, báo cáo khảo sát năm mã hóa AMO gồm Binomial, Binary, Sequential, Commander và Product, sử dụng Glucose3 thông qua PySAT. Trong phiên bản thực nghiệm hiện tại, Commander và Product đều được cài đặt theo cấu trúc đệ quy; Sequential sử dụng bộ đếm tuần tự và chuyển sang mã hóa theo cặp đối với nhóm không quá ba biến. Các cấu hình SAT được đối sánh với bốn bộ giải CPLEX CP Optimizer, CPLEX MIP, Gurobi và OR-Tools CP-SAT.

\subsection{Động cơ và câu hỏi nghiên cứu}

Binomial có cách xây dựng trực tiếp: mỗi cặp biến trong một nhóm AMO tạo ra một mệnh đề loại trừ. Cách biểu diễn này không cần biến phụ nhưng số mệnh đề tăng nhanh khi kích thước nhóm tăng. Các mã hóa còn lại tổ chức ràng buộc thông qua biến phụ, chuỗi trạng thái, nhóm hoặc tọa độ. Những thay đổi này có thể giảm kích thước CNF, đồng thời làm thay đổi quá trình lan truyền và tìm kiếm của bộ giải. Vì vậy, số mệnh đề là một tiêu chí cần đo, nhưng chưa đủ để kết luận một mã hóa có hiệu năng tốt hơn.

Từ đó, báo cáo tập trung vào bốn câu hỏi nghiên cứu:
\begin{enumerate}
    \item \textbf{Kích thước biểu diễn:} Số biến chính, biến phụ và mệnh đề CNF của từng mã hóa thay đổi như thế nào theo $N$? Các công thức đếm phản ánh cấu trúc cài đặt ra sao?
    \item \textbf{Chi phí xử lý:} Thời gian sinh CNF, nạp CNF và tìm nghiệm đóng góp như thế nào vào thời gian thực hiện? Việc giảm số mệnh đề có đi cùng với việc giảm thời gian giải trong các phép thử hay không?
    \item \textbf{Khả năng giải trong giới hạn tài nguyên:} Các cấu hình SAT và các mô hình CP/MIP khác nhau thế nào về số trường hợp hoàn tất, thời gian và bộ nhớ quan sát được dưới cùng giới hạn thực nghiệm?
    \item \textbf{Độ ổn định và khả năng mở rộng:} Kết quả giữa các lần lặp biến động ra sao, và với những lượt không hoàn tất, tiến trình đã đạt đến giai đoạn nào trước khi kết thúc?
\end{enumerate}

\subsection{Phạm vi và đóng góp của báo cáo}

Phạm vi nghiên cứu là tìm một cách đặt hậu hợp lệ trên bàn cờ trống cho mỗi kích thước được chọn. Mỗi cấu hình SAT sử dụng cùng các biến ô cờ và cùng các yêu cầu ALO/AMO của bài toán; khác biệt giữa chúng nằm ở cách mã hóa AMO. Đối với các bộ giải đối sánh, mô hình CP sử dụng biến vị trí và ràng buộc khác nhau từng đôi, còn mô hình MIP sử dụng biến nhị phân cùng các ràng buộc tuyến tính. Các mô hình cùng biểu diễn yêu cầu tìm nghiệm hợp lệ, nhưng có cấu trúc nội tại khác nhau.

Bộ dữ liệu mới gồm 18 kích thước bàn cờ từ $N=4$ đến $N=1331$, chín cấu hình và ba lần lặp cho mỗi cặp kích thước--cấu hình, tổng cộng 486 lượt chạy. Mỗi lượt được cấu hình ngân sách thời gian toàn lượt 100 giây và giới hạn bộ nhớ 6144 MiB. Chương trình không loại trước một trường hợp dựa trên số mệnh đề CNF. Commander và Product trong toàn bộ dữ liệu này là các bản đệ quy; do đó, kết quả được dùng để so sánh các cài đặt hiện tại, không phải để định lượng riêng tác động của việc chuyển từ một tầng sang đệ quy.

Đóng góp của báo cáo gồm ba nội dung. Thứ nhất, báo cáo trình bày đầy đủ cách chuyển N-Queens thành SAT, giải thích vai trò của ALO, AMO và năm cách mã hóa được sử dụng. Thứ hai, báo cáo xây dựng cách đọc kết quả nhất quán: phân biệt kích thước CNF đầy đủ theo công thức với tiến độ sinh/nạp thực tế, đồng thời phân biệt thời gian xử lý mô hình và giải với thời gian toàn lượt. Thứ ba, báo cáo phân tích dữ liệu thực nghiệm theo thời gian, kích thước biểu diễn, bộ nhớ và mức độ hoàn tất; các nghiệm tìm được được kiểm tra lại bằng điều kiện hợp lệ của N-Queens.

\subsection{Tài liệu liên quan và tổ chức báo cáo}
Tài liệu N-Queens của Trần Ngọc Bách \citep{bach2024} là nguồn tham khảo cho cách lựa chọn kích thước bàn cờ và trình bày bảng so sánh encoding. Slide ALO/AMO của học phần \citep{amoSlides} cung cấp các nguyên lý mã hóa nền tảng; các lựa chọn cụ thể của cài đặt hiện tại được nêu trong Methodology. Mẫu bài báo môn học \citep{courseTemplate} định hướng việc tách mô hình, thiết lập thực nghiệm, kết quả và Data Availability. Các thời gian của tài liệu tham khảo không được gộp vào dữ liệu đo trên máy của người thực hiện.


Phần~\ref{sec:background} (Background) trình bày nền tảng về N-Queens, SAT, CNF, ALO/AMO và các họ bộ giải. Phần~\ref{sec:methodology} (Methodology) xây dựng mô hình SAT, mô tả chi tiết năm mã hóa và các mô hình CP/MIP đối sánh. Phần~\ref{sec:experiments} (Experiments and Results) trình bày thiết lập thực nghiệm, cách đo, các bảng kết quả và phân tích so sánh. Phần~\ref{sec:conclusion} (Conclusion) tổng hợp những kết luận được hỗ trợ bởi dữ liệu và nêu hướng phát triển tiếp theo.
